% Plantilla LaTeX canónica de CV para el generador Job-up.
% Los marcadores se sustituyen literalmente desde datos-generacion.json.
\documentclass[11pt,a4paper]{article}
\usepackage[utf8]{inputenc}
\usepackage[T1]{fontenc}
\usepackage[margin=1.6cm]{geometry}
\usepackage{xcolor}
\usepackage{enumitem}
\usepackage{hyperref}
\setlength{\parindent}{0pt}
\setlength{\parskip}{4pt}
\definecolor{primary}{HTML}{1F2937}
\definecolor{secondary}{HTML}{5B6573}
\begin{document}

{\color{primary}\LARGE\textbf{[NOMBRE]}}\\
{\color{secondary}\large [TITULAR]}\\
[EMAIL] \textbar{} [TELÉFONO] \textbar{} [LINKEDIN]

\section*{Perfil profesional}
[PERFIL PROFESIONAL]

\section*{Propuesta de valor}
[PROPUESTA DE VALOR]

\section*{Experiencia profesional relevante}
\textbf{[EXPERIENCIA 1 CABECERA]} -- [EXPERIENCIA 1 DESCRIPCION]

\textbf{[EXPERIENCIA 2 CABECERA]} -- [EXPERIENCIA 2 DESCRIPCION]

\textbf{[EXPERIENCIA 3 CABECERA]} -- [EXPERIENCIA 3 DESCRIPCION]

\textbf{[EXPERIENCIA 4 CABECERA]} -- [EXPERIENCIA 4 DESCRIPCION]

\textbf{[EXPERIENCIA 5 CABECERA]} -- [EXPERIENCIA 5 DESCRIPCION]

\textbf{[EXPERIENCIA 6 CABECERA]} -- [EXPERIENCIA 6 DESCRIPCION]

\section*{Competencias y herramientas}
[COMPETENCIA 1]\\
[COMPETENCIA 2]\\
[COMPETENCIA 3]\\
[COMPETENCIA 4]

\section*{Formación}
[FORMACION 1]\\
[FORMACION 2]\\
[FORMACION 3]

[INFORMACION ADICIONAL]

\end{document}
